% =====================================================================
%  capitulos/cap12.tex — Capítulo 12. Las relaciones interoracionales:
%  coordinación
%  Autora: Natalia Giollo
%  Fuente: Capitulo_12_normalizado.docx
% =====================================================================
\capitulo{Las relaciones interoracionales: coordinación}{Natalia Giollo}

\section{Introducción}

En este capítulo comenzamos el estudio de las relaciones que vinculan distintas
oraciones, también conocidas como \termino{relaciones interoracionales}. Para
ello, partimos de mostrar cómo una predicación simple puede ir ampliándose para
conformar estructuras más complejas, que incluyan otra u otras predicaciones, ya
sea secundarias, o incluidas, como en la subordinación,\footnote{La subordinación
se aborda en el \capref{13}.} o conservando cada una su independencia
sintáctica, como en la coordinación. Luego, nos centramos en la coordinación y
abordamos las particularidades de los distintos subtipos (copulativa, disyuntiva
y adversativa). Asimismo, también introducimos el concepto de marcador discursivo
y nos ocupamos de la identificación de las instancias de elipsis asociadas a la
coordinación.

\section{El pasaje de la oración simple a la oración compuesta}\label{sec:12-pasaje}

La oración, tal como se la definió en el \capref{6}, §\ref{sec:06-definicion},
constituye una unidad sintáctico-semántica, en la que un SN, como \forma{mi
hermano} en \exref{ex:12-hermano1}, y un SV, como \forma{lavó los platos}
\exref{ex:12-hermano1}, establecen una relación predicativa articulada a través
de las categorías funcionales de Tiempo y Concordancia para conformar un
sintagma de flexión (SF): \forma{Mi hermano lavó los platos}.

\begin{exe}
  \ex\label{ex:12-hermano1} {}[\etq{SF/O} [\etq{SN/Suj} Mi hermano]
      [\etq{SV/Pred} lavó los platos]]
\end{exe}

Ahora bien, además de expresar una predicación simple, la oración también puede
incorporar una segunda instancia predicativa, como cuando contiene algún
predicativo no obligatorio, que aporta una predicación adicional o secundaria
\exref{ex:12-hermano2} (véase el \capref{9}, §\ref{sec:09-primaria}):

\begin{exe}
  \ex\label{ex:12-hermano2} {}[\etq{SF/O} [\etq{SN/Suj} Mi hermano]
      [\etq{SV/Pred} lavó los platos [\etq{SA/PSNO} contento]]]
\end{exe}

En \exref{ex:12-hermano1} el SF contiene una única predicación, organizada a
partir de \forma{lavar}, un verbo transitivo que expresa una acción realizada
por un sujeto agente ---\forma{mi hermano}--- sobre un tema (el SN/COD \forma{los
platos}). En \exref{ex:12-hermano2} se mantiene la predicación de
\exref{ex:12-hermano1}, pero aparece una segunda instancia predicativa: la
\termino{predicación secundaria}, formada por el adjetivo \forma{contento},
referido al mismo sujeto de \exref{ex:12-hermano1}. Según vimos en el
\capref{9}, §\ref{sec:09-sintaxis}, el predicativo forma una oración reducida,
articulada sin flexión, ya que la oración reducida carece de la manifestación de
las categorías de tiempo, aspecto y modo. Asimismo, el sujeto en esta segunda
instancia predicativa ---\forma{(estar) contento}--- tiene otro papel temático:
es un \textsc{paciente}.

Si bien en \exref{ex:12-hermano2} hay dos predicaciones, el SF se sigue
considerando una oración simple, debido a que la segunda predicación no es
autónoma o completa, puesto que carece de flexión y se temporaliza a partir del
tiempo de la predicación principal. En efecto, en \exref{ex:12-hermano2}, la
única interpretación posible del evento \forma{(estar) contento} es que se
encuentra en el mismo tiempo pasado establecido por el evento de \forma{lavar}.
En este sentido, \exref{ex:12-hermano2} se podría parafrasear como se muestra en
\exref{ex:12-hermano3}.

\begin{exe}
  \ex\label{ex:12-hermano3} {}[\etq{SF/O} Mi hermano lavó los platos].
      [\etq{SF/O} (\textit{pro}) Estaba contento].
\end{exe}

En \exref{ex:12-hermano3}, hay dos sintagmas de flexión independientes y cada uno
constituye una oración simple. La estructura argumental que manifiestan es la
siguiente: \forma{lavar} ⟨alguien⟩ ⟨algo⟩ y \forma{estar contento} ⟨alguien⟩. En
la primera oración, los argumentos del verbo son ⟨\forma{mi hermano}⟩ y
⟨\forma{los platos}⟩. En la segunda, como ya adelantamos, el argumento
correspondiente al sujeto está representado sintácticamente por un sujeto tácito
o desinencial que, de acuerdo con lo desarrollado en el \capref{6},
§\ref{sec:06-tacito}, se corresponde con una forma pronominal abstracta, un
\textit{pro}, es decir, un pronombre fonéticamente vacío, pero semánticamente
interpretable y recuperable a partir del contexto de situación o de lo dicho
anteriormente. En este caso, \textit{pro} recoge el valor referencial del SN
\forma{mi hermano}, el sujeto de la oración inmediatamente anterior. A partir de
esa referencia entendemos que estas oraciones están vinculadas: lo que era el
contenido de una única oración simple \exref{ex:12-hermano2} ahora aparece
distribuido en dos oraciones, simples también, una contigua a la otra, en una
relación de \termino{yuxtaposición}.

\begin{destacado}
\textbf{Para recordar.} \textcite[334]{bosque2019} define la yuxtaposición como

\begin{cita}
Unión sintáctica de dos o más elementos en la cual se omite la marca formal (una
conjunción o una preposición), que expresa la relación gramatical que se
establece entre ellos [\ldots].
\end{cita}

Los elementos que se unen mediante yuxtaposición pueden ser sintagmas nominales,
adjetivales, preposicionales, adverbiales y verbales, además de oracionales. Al
igual que en la coordinación, los segmentos yuxtapuestos son elementos contiguos
del mismo nivel. La relación semántica que se da entre ellos es también semejante
a la que se establece en la coordinación (adición, exclusión, etc.).
\end{destacado}

Siguiendo esta línea de razonamiento, podríamos parafrasear
\exref{ex:12-hermano3} con una estructura como la presentada en
\exref{ex:12-hermano4}, en la que se emplea una conjunción para explicitar la
vinculación entre ambas oraciones:

\begin{exe}
  \ex\label{ex:12-hermano4} Mi hermano lavó los platos y estaba contento.
\end{exe}

En \exref{ex:12-hermano4}, a partir de la definición de oración que hemos
adoptado, se puede considerar que tenemos una \termino{oración compuesta} (OC)
constituida por dos oraciones simples (OS) coordinadas entre sí, cada una con su
sujeto: el SN \forma{mi hermano}, para la primera oración, y el pronombre vacío
\textit{pro}, para la segunda. Es cierto que, en términos referenciales, ambas
formas denotan la misma entidad del mundo extralingüístico y que comparten,
necesariamente, los rasgos de persona y número (tercera persona del singular)
responsables de la concordancia con los verbos \forma{lavó} y \forma{estaba},
respectivamente. Sin embargo, no se trata del mismo sujeto en términos
semánticos, porque cada instancia aporta un papel temático distinto. En este
sentido, en [\forma{mi hermano lavó los platos}] el sujeto es la entidad que
realiza la acción, el agente de \forma{lavar}, mientras que en [\forma{pro
estaba contento}], el sujeto es la entidad a la que se le atribuye un estado, con
el papel paciente de \forma{estar contento}. Es importante destacar que, tal como
se ha desarrollado en el \capref{7}, cada predicado tiene una estructura
argumental y temática propia, independientemente de que sus respectivos sujetos
compartan (o no) valores referenciales, sintácticos y semánticos.

Por último, también es posible reformular \exref{ex:12-hermano4} incluyendo una
oración en otra, es decir, utilizando una estructura incluida o subordinada a la
principal, como se muestra en \exref{ex:12-hermano5}:

\begin{exe}
  \ex\label{ex:12-hermano5}
    \begin{xlist}
      \ex\label{ex:12-hermano5-a} Mi hermano, [que estaba contento], lavó los
          platos.
      \ex\label{ex:12-hermano5-b} Mi hermano, [que lavó los platos], estaba
          contento.
      \ex\label{ex:12-hermano5-c} [Cuando lavó los platos] mi hermano estaba
          contento.
    \end{xlist}
\end{exe}

En los ejemplos de arriba se encuentran las mismas dos predicaciones
(\forma{lavar} ⟨alguien⟩ ⟨algo⟩, por un lado, y \forma{estar contento}
⟨alguien⟩, por el otro), ambas realizadas como predicados independientes. Pero
ahora una de ellas es una oración subordinada (que aparece entre corchetes), es
decir, una oración incluida en otra principal de la que forma parte
constituyendo una estructura mayor: una oración compuesta, en la que hay una
oración principal (predicación principal) y dentro de esta una oración
subordinada (predicación secundaria). Como se desprende del contraste entre
\exref{ex:12-hermano5-a}, \exref{ex:12-hermano5-b} y \exref{ex:12-hermano5-c},
el hablante puede crear distintas configuraciones sintácticas, preservando ---en
líneas generales--- el valor de cada predicación.

\begin{recuadro}[Sobre el concepto de oración compuesta]
Si bien en las gramáticas tradicionales se distinguía entre oración compuesta,
formada por oraciones coordinadas, y oración compleja, constituida por una
oración principal y una o varias subordinadas, en este libro seguimos los
planteos actuales de la \textit{NGLE} \parencite[§\,1.13ñ]{rae2025}, que adopta el
término \textit{oración compuesta} para todos los casos de combinación de
oraciones. Igualmente, en \textcite{bosque2019} se plantea que las oraciones que
forman una oración compuesta pueden estar vinculadas por coordinación
(\forma{María estudia y su hermana trabaja; He visto la película, pero no me ha
gustado}) o por subordinación (\forma{Me prometiste que vendrías; El contrato
que firmó no era válido}). Las oraciones compuestas también se pueden vincular
por yuxtaposición (\forma{Tu padre está cansado; pasa a buscarlo con el
coche}).
\end{recuadro}

En resumen, lo que se muestra de \exref{ex:12-hermano1} a \exref{ex:12-hermano5}
es que una oración simple, como \exref{ex:12-hermano1} \forma{Mi hermano lavó los
platos}, se puede ir transformando en estructuras que incorporen otras
predicaciones, en grados de complejidad creciente. Por ejemplo, en
\exref{ex:12-hermano2} \forma{Mi hermano lavó los platos contento}, la
predicación se expande con la adición del predicativo, que se añade a la
predicación principal. Asimismo, ambas predicaciones pueden ser recogidas en dos
oraciones independientes entre las que medie una mera yuxtaposición
\exref{ex:12-hermano3}: \forma{Mi hermano lavó los platos. Estaba contento}. O
bien, como hemos visto en \exref{ex:12-hermano4} y \exref{ex:12-hermano5},
podemos vincular ambas predicaciones, ya sea mediante conjunciones coordinantes
o subordinantes: la conjunción \forma{y}, que coordina dos predicaciones
principales, en \exref{ex:12-hermano4}, y los relativos \forma{que} y
\forma{cuando}, que establecen la relación entre la oración subordinada y la
principal, en \exref{ex:12-hermano5}. Si bien las predicaciones son las mismas,
lo que va variando es la manera en que se articulan entre sí, de modo que se
están relacionando oraciones (cláusulas o también proposiciones, según las
denominaciones dadas en las distintas tradiciones gramaticales) mediante
diferentes relaciones interoracionales.

En los siguientes apartados, nos centramos en analizar el complejo campo de las
relaciones interoracionales que se establecen mediante el mecanismo de la
coordinación, con énfasis en el análisis de los distintos tipos de conjunciones
coordinantes. Queda para el próximo capítulo el tratamiento más extenso de los
distintos tipos de oraciones subordinadas.

\begin{ejercicio}\label{ej:12-predicaciones}
En las siguientes oraciones, distinga la predicación primaria y la secundaria.
Luego, proponga una versión por yuxtaposición, otra por coordinación y otra por
subordinación, según se muestra en el ejemplo resuelto~1.

\begin{exe}
  \ex Los niños caminaron hacia el parque felices.
  \ex Los transeúntes miraban la escena asombrados.
  \ex María colgó el cuadro bien centrado.
\end{exe}

\begin{resuelto}
1. Predicación primaria: \forma{Los niños caminaron hacia el parque}.
Predicación secundaria: \forma{Los niños estaban felices}. Yuxtaposición:
\forma{Los niños caminaron hacia el parque, estaban felices}. Coordinación:
\forma{Los niños caminaron hacia el parque y estaban felices}. Subordinación: a)
\forma{Los niños, que estaban felices, caminaron hacia el parque}; b) \forma{Los
niños, que caminaron hacia el parque, estaban felices}; c) \forma{Cuando
caminaron hacia el parque, los niños estaban felices}.
\end{resuelto}
\end{ejercicio}

\section{La coordinación}\label{sec:12-coordinacion}

La \termino{coordinación} puede definirse como «el proceso gramatical por el cual
dos o más constituyentes sintácticos (los "coordinandos") se unen con una
conjunción formando una frase compleja» \parencite[301]{camacho2016}. Mediante la
coordinación se pueden vincular distintos tipos de constituyentes, que
generalmente son sintagmas de la misma categoría, de modo tal que la
coordinación se caracteriza como \termino{homocategorial}.

\begin{exe}
  \ex\label{ex:12-juan-pedro} {}[\etq{SSNN} Juan \textit{y} Pedro] fueron al río.
  \ex\label{ex:12-rio} El río estaba [\etq{SSAA} calmo, \textit{pero} muy frío].
  \ex\label{ex:12-ahora} Lo reclamamos [\etq{SSAdv} ahora \textit{o} nunca].
  \ex\label{ex:12-animo} Juan se encuentra [\etq{SSPP} de buen ánimo \textit{y} en
      buen estado de salud].
  \ex\label{ex:12-maria} {}[[\etq{SF} María llegó a la oficina] \textit{y}
      [\etq{SF} encendió la computadora]].
\end{exe}

La conjunción no solo coordina constituyentes que pertenecen a la misma clase,
sino que también une constituyentes de distinta clase, pero con una función
equivalente, por lo que también existe la coordinación
\termino{heterocategorial}.

\begin{exe}
  \ex\label{ex:12-cansado} Entró [[\etq{SA} cansado] \textit{y} [\etq{SP} sin
      prisa]]. $\rightarrow$ Coordinación de un SA y un SP.
  \ex\label{ex:12-rapido} Llegó [[\etq{SAdv} rápidamente] \textit{y} [\etq{SP}
      con buena disposición]]. $\rightarrow$ Coordinación de un SAdv y un SP.
  \ex\label{ex:12-bebida} Compramos [\etq{SN} la bebida] \textit{y} [\etq{SF}
      todo lo que estaba en oferta]. $\rightarrow$ Coordinación de un SN y un
      SF.
\end{exe}

\subsection{La estructura de la coordinación}\label{sec:12-estructura}

En relación con la coordinación, tanto la gramática tradicional como la
estructural planteaban una estructura plana en la que los constituyentes unidos
por la conjunción estaban en el mismo nivel y no había un núcleo que diera la
categoría al sintagma resultante de dicha unión:

\begin{exe}
  \ex {}[[\etq{SN} Los estudiantes] \textit{y} [\etq{SN} sus profesores]]
      cantaron el himno nacional en la ceremonia.
  \ex {}[[\etq{SF} El gobierno aumentó los impuestos], \textit{pero} [\etq{SF}
      la ciudadanía se manifestó en contra]].
\end{exe}

Sin embargo, una segmentación como la de arriba no se adecua a los planteos de la
teoría sintagmática propuesta por la gramática generativa (y presentada en el
\capref{5}, §\ref{sec:05-xbarra}), ya que no permite identificar un núcleo que
otorgue la categoría a toda la construcción. Al respecto, recordemos que, según
el esquema de la teoría X\prima, todos los sintagmas son endocéntricos, es decir,
tienen un núcleo, que puede estar solo o ir acompañado de un complemento en el
nivel X\prima, y que ambos, núcleo y complemento, pueden recibir un especificador
que incide sobre ellos, tal como se muestra en \exref{ex:12-xbarra}.

\begin{exe}
  \ex\label{ex:12-xbarra}
      \begin{forest} gramatica
        [SX
          [Esp]
          [X\prima
            [X\\(núcleo)]
            [Comp]]]
      \end{forest}
\end{exe}

Siguiendo, entonces, estos planteos, necesitamos recurrir a un nuevo sintagma de
naturaleza funcional: el \termino{sintagma de conjunción} (SConj). Su núcleo es
una palabra funcional, la \termino{conjunción coordinante} (Conj), que toma como
complemento al segundo constituyente de la coordinación y, como especificador,
al primer constituyente, según el esquema de \exref{ex:12-sconj}.

\begin{exe}
  \ex\label{ex:12-sconj}
      \begin{forest} gramatica
        [SConj
          [Esp\\(1.\textsuperscript{er} coordinando)]
          [Conj\prima
            [Conj\\(núcleo)]
            [Comp\\(2.\textsuperscript{o} coordinando)]]]
      \end{forest}
\end{exe}

Es importante resaltar que el SConj se proyecta en toda instancia en que aparece
una conjunción, más allá de la naturaleza de los elementos coordinados. En todos
los casos, en términos estructurales, el elemento que aparece a la derecha de la
conjunción coordinante ocupa la posición de complemento, mientras que el
elemento a la izquierda se ubica en la posición de especificador. Veamos, en
\exref{ex:12-an-juan}--\exref{ex:12-an-maria}, cómo se aplica este análisis a las
oraciones de \exref{ex:12-juan-pedro} a \exref{ex:12-maria}.

\begin{exe}
  \ex\label{ex:12-an-juan} {}[\etq{SConj} [\etq{SN/Esp} Juan] \textit{y}\etq{Conj/n}
      [\etq{SN/Comp} Pedro]] fueron al río.
  \ex\label{ex:12-an-rio} El río estaba [\etq{SConj} [\etq{SA/Esp} calmo],
      \textit{pero}\etq{Conj/n} [\etq{SA/Comp} muy frío]].
  \ex Lo reclamamos [\etq{SConj} [\etq{SAdv/Esp} ahora] \textit{o}\etq{Conj/n}
      [\etq{SAdv/Comp} nunca]].
  \ex Juan se encuentra [\etq{SConj} [\etq{SP/Esp} de buen ánimo]
      \textit{y}\etq{Conj/n} [\etq{SP/Comp} en buen estado de salud]].
  \ex\label{ex:12-an-maria} {}[\etq{SConj} [\etq{SF/Esp} María llegó a la
      oficina] \textit{y}\etq{Conj/n} [\etq{SF/Comp} encendió la computadora]].
\end{exe}

Ahora bien, analizar dos (o más) constituyentes como los de la coordinación, con
la misma o equivalente naturaleza categorial, como una estructura jerárquica con
diferentes niveles puede, en principio, parecer contraintuitivo. Por eso,
dedicamos el próximo apartado a fundamentar estos planteos.

\subsection{Sobre la estructura jerárquica de la coordinación}\label{sec:12-jerarquica}

El mecanismo de la coordinación tiene la particularidad de presentar al mismo
tiempo características de equivalencia categorial entre sus miembros, pero
también, desde la perspectiva estructural o de conformación sintagmática, se
reconocen relaciones jerárquicas entre los coordinandos. En principio, aunque los
coordinandos compartan la misma naturaleza categorial, a partir de la estructura
propuesta en \exref{ex:12-sconj}, asumimos una relación jerárquica o asimétrica
entre ellos, en tanto el complemento y el especificador no se encuentran en el
mismo nivel, sino que ocupan diferentes posiciones estructurales.

Asimismo, la vinculación entre ambas oraciones coordinadas requiere una operación
sintáctica básica: el \termino{ensamble} o
\termino{fusión}.\pendiente{El original dice que el ensamble «ya [fue]
mencionado en el capítulo 6», pero no aparece allí} Como se explica en
\textcite[219]{bosque2009}, el ensamble es «el proceso fundamental por el que
formamos nuevos constituyentes». Cuando dos expresiones se ensamblan, se produce
un nuevo constituyente que hereda los rasgos categoriales de uno de los
elementos fusionados. En el caso de las estructuras coordinadas, primero se
ensambla la conjunción coordinante, que actúa como núcleo de la construcción, con
su complemento, que es el segundo miembro de la coordinación. Luego, el
constituyente resultante (Conj\prima) se ensambla con el primer miembro de la
coordinación, que está en la posición de especificador, dando lugar a la
proyección máxima del sintagma (SConj). En \exref{ex:12-arbol-maria} presentamos
la estructura arbórea de \exref{ex:12-maria}.

\begin{exe}
  \ex\label{ex:12-arbol-maria}
      \begin{forest} gramatica
        [SConj
          [SF\\Esp [\textit{María llegó a la oficina}, roof]]
          [Conj\prima
            [Conj\\\textit{y}]
            [SF\\Comp [\textit{encendió la computadora}, roof]]]]
      \end{forest}
\end{exe}

Ahora bien, asumir que el segundo coordinando es el complemento de la conjunción
implica plantear que entre este y la conjunción existe una relación más estrecha
---de complementación--- que la que se da entre la conjunción y el primer
coordinando. ¿Cómo podemos fundamentar esta diferencia estructural? Con el fin de
recoger evidencia, prestemos atención a los datos que nos puede proporcionar un
simple intercambio entre dos hablantes, A y B. En \exref{ex:12-dialogo-a}, el
hablante A hace una afirmación y su interlocutor B la continúa en una típica
réplica \exref{ex:12-dialogo-b}, que puede ser tanto enunciativa como
interrogativa y está conformada solo por la conjunción y su complemento. Según
vemos en la réplica de \exref{ex:12-dialogo-b}, la estructura formada por la
conjunción seguida del segundo coordinando (su complemento), pero sin incluir al
primer constituyente (que ya fue expresado por el primer emisor), resulta
perfectamente gramatical. En cambio, si la réplica del interlocutor fuera como
\exref{ex:12-dialogo-c} y estuviera formada por el primer constituyente seguido
de la conjunción (sin el complemento), no resultaría gramatical.

\begin{exe}
  \ex\label{ex:12-dialogo}
    \begin{xlist}
      \ex\label{ex:12-dialogo-a} A: ---Mañana vendrá María.
      \ex\label{ex:12-dialogo-b} B: ---\textit{Y} Pedro. / ---¿\textit{Y} Pedro?
      \ex[*]{\label{ex:12-dialogo-c} B\prima: ---Juan \textit{y}.}
    \end{xlist}
\end{exe}

Veamos otros dos ejemplos, que nos permitirán reforzar la argumentación
anterior:

\begin{exe}
  \ex\label{ex:12-juguete}
    \begin{xlist}
      \ex Todavía no sabemos si fue un accidente, \textit{pero} rompimos el
          juguete en mil pedazos. (H.~Casciari, \textit{Teníamos un juguete})
      \ex[*]{\label{ex:12-juguete-b} Todavía no sabemos si fue un accidente,
          \textit{pero}.}
      \ex\label{ex:12-juguete-c} \textit{Pero} rompimos el juguete en mil
          pedazos.
    \end{xlist}
  \ex\label{ex:12-padres}
    \begin{xlist}
      \ex Algunos padres todavía se amontonan frente a las puertas \textit{y}
          gritan los nombres de sus hijos. (S.~Schweblin, \textit{Mariposas})
      \ex[*]{\label{ex:12-padres-b} Algunos padres todavía se amontonan frente a
          las puertas \textit{y}.}
      \ex\label{ex:12-padres-c} \textit{Y} gritan los nombres de sus hijos.
    \end{xlist}
\end{exe}

A partir del contraste de gramaticalidad entre los casos (b) y (c) de los
ejemplos \exref{ex:12-juguete} y \exref{ex:12-padres}, podemos nuevamente
comprobar que la conjunción coordinante establece una relación muy estrecha con
su complemento y que solamente es gramatical la estructura en que va seguida de
su complemento, con ausencia de especificador, como en \exref{ex:12-juguete-c} y
\exref{ex:12-padres-c}; mientras que ni \exref{ex:12-juguete-b} ni
\exref{ex:12-padres-b} resultan gramaticales debido, justamente, a la ausencia
del complemento de la conjunción, que no puede ser omitido.\footnote{El
complemento puede omitirse si queremos generar suspenso: \forma{Algunos padres
gritan los nombres de sus hijos y\ldots}, caso en que se considera que el
complemento queda implícito.}

Por otra parte, también debemos señalar que existe una equivalencia funcional
entre los coordinandos cuando se toman por separado, es decir que ambos podrían
cumplir independientemente la misma función sintáctica. En \exref{ex:12-rio},
reproducida aquí como \exref{ex:12-equiv-a}, tanto el SA \forma{calmo} como el
SA \forma{muy frío} podrían funcionar como PSO en un contexto sin coordinación. A
su vez, en \exref{ex:12-equiv-b}, tanto el SA \forma{cansado} como el SP
\forma{sin prisa} podrían actuar como PSNO: \forma{Entró cansado; Entró sin
prisa}.

\begin{exe}
  \ex
    \begin{xlist}
      \ex\label{ex:12-equiv-a} El río estaba [\etq{SA/PSO} calmo]. El río estaba
          [\etq{SA/PSO} muy frío].
      \ex\label{ex:12-equiv-b} Entró [\etq{SA/PSNO} cansado]. Entró
          [\etq{SP/PSNO} sin prisa].
    \end{xlist}
\end{exe}

Ahora bien, cuando los elementos están coordinados, asumimos que el SConj es el
que cumple la función sintáctica que ambos tenían por separado. Al mismo tiempo,
al establecer entre ambos miembros una relación jerárquica dentro del SConj, cuyo
núcleo es la conjunción, el primer coordinando pasa a funcionar como Esp y el
segundo, como Comp:

\begin{exe}
  \ex
    \begin{xlist}
      \ex El río estaba [\etq{SConj/PSO} [\etq{SA/Esp} calmo],
          \textit{pero}\etq{Conj/n} [\etq{SA/Comp} muy frío]].
      \ex Entró [\etq{SConj/PSNO} [\etq{SA/Esp} cansado] \textit{y}\etq{Conj/n}
          [\etq{SP/Comp} sin prisa]].
    \end{xlist}
\end{exe}

\begin{ejercicio}\label{ej:12-desantis}
En las siguientes oraciones de \textit{La sexta lámpara}, de Pablo De Santis,
señale entre corchetes los dos miembros de cada coordinación y la estructura
resultante, como se muestra en el ejemplo resuelto~1.

\begin{exe}
  \ex El viejo Treviso supo aprovechar nuestro fanatismo y empezó a pedirnos
      pequeños préstamos a cambio de una información siempre escatimada y a
      menudo contradictoria.
  \ex Hicimos una colecta para pagarle el legado, sin saber si comprábamos un
      tesoro o un montón de diarios viejos.
  \ex Habíamos subido el baúl por la escalera casi a los saltos, pero la visión
      de aquellos papeles nos paralizó.
  \ex Balestri sigue siendo un desconocido, y de su torre no poseemos más que la
      sombra de una sombra.
\end{exe}

\begin{resuelto}
1. [\etq{SConj/OC} [\etq{SF/Esp} El viejo Treviso supo aprovechar nuestro
fanatismo] y\etq{Conj/n} [\etq{SF/Comp} empezó a pedirnos pequeños préstamos a
cambio de una información [\etq{SConj/M} [\etq{SV(participio)/Esp} siempre
escatimada] y\etq{Conj/n} [\etq{SA/Comp} a menudo contradictoria]]]].
\end{resuelto}
\end{ejercicio}

\begin{ejercicio}\label{ej:12-enriquez}
Señale entre corchetes las instancias de coordinación del siguiente fragmento de
\textit{Éste es el mar}, de Mariana Enriquez, tal como se observa en el ejemplo
resuelto.

\begin{textolectura}
Tenían una madre que rara vez era vista o convocada; la madre Hécate. En el
Enjambre se creía que era una Leyenda, una mentira, pero las Luminosas se reían
de tanta ignorancia. Y tenían hermanas mayores: Perséfone bajo la tierra, con
flores en el pelo; y las Imago, las que se escondían en lugares secos y oscuros.
\end{textolectura}

\begin{resuelto}
[\etq{Oración 1} Tenían una madre que rara vez era
[\etq{SConj}\footnote{En este caso, la conjunción une dos formas verbales: los
dos auxiliados de la perífrasis de voz pasiva \forma{era vista o convocada}.
Como la coordinación se da dentro de la perífrasis verbal, el SConj no tiene
función sintáctica. Las perífrasis verbales se tratan en el
\capref{11}.\pendiente{El original dice «los dos auxiliares» y etiqueta
\textit{vista} y \textit{convocada} como Vaux}} [\etq{V/Esp} vista]
o\etq{Conj/n} [\etq{V/Comp} convocada]]; la madre Hécate].
\end{resuelto}
\end{ejercicio}

\begin{destacado}
\textbf{En síntesis.} Estructuralmente, la coordinación es una relación que une
dos constituyentes (generalmente, de la misma categoría) entre los que establece
una relación jerárquica. Por tanto, podemos afirmar que, al igual que el resto de
los sintagmas estudiados en otros capítulos, la coordinación se adecua al esquema
de X\prima, respetando el principio de endocentricidad: posee un núcleo funcional
---la conjunción--- que toma un complemento obligatorio (el segundo coordinando) y
un especificador que puede, eventualmente, omitirse (el primer coordinando).
\end{destacado}

\subsection{Las conjunciones coordinantes}\label{sec:12-conjunciones}

Para analizar la relación que se establece entre los coordinandos es fundamental
identificar la naturaleza del núcleo de la coordinación, es decir, de las
distintas conjunciones coordinantes. En principio, las conjunciones coordinantes
pueden ser de dos tipos:

\begin{itemize}
  \item \termino{conjunciones coordinantes simples}: \forma{y}, con su variante
        \forma{e}; \forma{ni} (= \forma{y no}); \forma{o}, con su variante
        \forma{u}; \forma{pero}; \forma{sino}; \forma{mas};
  \item \termino{conjunciones coordinantes correlativas}: \forma{ni\ldots\
        ni\ldots; tanto\ldots\ como\ldots; tanto\ldots\ cuanto\ldots;
        así\ldots\ como\ldots; sea\ldots\ sea\ldots; ya\ldots\ ya\ldots;
        ora\ldots\ ora\ldots; bien\ldots\ bien\ldots; o\ldots\ o\ldots}
\end{itemize}

Es importante destacar que, más allá del hecho de que las conjunciones
coordinantes sean simples o compuestas, la coordinación puede ser

\begin{itemize}
  \item \termino{binaria}, cuando se establece entre solo dos coordinandos;
  \item \termino{múltiple}, cuando vincula varios coordinandos.
\end{itemize}

Algunas conjunciones pueden usarse de forma binaria o múltiple. La conjunción
\forma{y}, por ejemplo, es simple, pero puede participar tanto en coordinaciones
binarias, como las de \exref{ex:12-ajedrez-a} y \exref{ex:12-ajedrez-b}, como en
múltiples, como las de \exref{ex:12-ajedrez-c} y \exref{ex:12-ajedrez-d}:

\begin{exe}
  \ex
    \begin{xlist}
      \ex\label{ex:12-ajedrez-a} Juan juega al ajedrez \textit{y} a las damas.
      \ex\label{ex:12-ajedrez-b} Juan juega al ajedrez \textit{y} Pedro juega a
          las damas.
      \ex\label{ex:12-ajedrez-c} Juan juega al ajedrez, las damas, el backgammon
          \textit{y} el póker.
      \ex\label{ex:12-ajedrez-d} Juan juega al ajedrez, Pedro prefiere las damas
          \textit{y} su hermana María ama el póker.
    \end{xlist}
\end{exe}

Las conjunciones que pueden aparecer en coordinaciones múltiples son: \forma{y,
ni, tanto\ldots\ como, o, o sea, bien\ldots\ bien}. En cambio, las conjunciones
\forma{pero} y \forma{sino} solo pueden participar en coordinaciones binarias,
como se aprecia en los ejemplos de \exref{ex:12-binarias}:\pendiente{\textit{o
sea} figura entre las conjunciones, pero en la tabla~\ref{tab:12-marcadores} se
clasifica como marcador explicativo}

\begin{exe}
  \ex\label{ex:12-binarias}
    \begin{xlist}
      \ex Juan juega al ajedrez, \textit{pero} Pedro prefiere jugar a las damas.
      \ex[*]{Juan juega al ajedrez, \textit{pero} Pedro prefiere jugar a las
          damas, \textit{pero} su hermana María ama jugar al póker.}
      \ex Juan no juega al ajedrez, \textit{sino} a las damas.
      \ex[*]{Juan no juega al ajedrez, \textit{sino} a las damas, \textit{sino}
          al póker.}
    \end{xlist}
\end{exe}

Por otra parte, según su funcionamiento sintáctico-semántico, las conjunciones se
clasifican en tres grupos:

\begin{itemize}
  \item \termino{copulativas}: \forma{y, e, ni, tanto\ldots\ como\ldots};
  \item \termino{disyuntivas}: \forma{o, u, bien\ldots\ bien\ldots,
        sea\ldots\ sea\ldots};
  \item \termino{adversativas}: \forma{pero, sino, mas}.
\end{itemize}

Cuando la conjunción coordinante es copulativa, la relación de coordinación se
interpreta como una suma o adición. En cambio, cuando es disyuntiva, se entiende
que hay una elección o alternancia. En el caso de las conjunciones adversativas,
se plantea una relación de oposición entre los coordinandos. En los próximos
apartados, describimos los distintos tipos de conjunciones coordinantes y las
características de las estructuras de coordinación que conforman.

\subsubsection{La coordinación copulativa}\label{sec:12-copulativa}

En la coordinación copulativa se establece entre los coordinandos un vínculo de
adición. La conjunción coordinante copulativa simple \forma{y} constituye el
recurso más empleado para la coordinación. Otras conjunciones copulativas son:
\forma{ni; tanto\ldots\ como\ldots};\footnote{\textcite[194]{jimenezjulia2023}
explica que el valor de \forma{tanto\ldots\ como\ldots} se solapa con el de la
conjunción copulativa \forma{y}, en tanto ambas formas permiten combinar acciones
que se producen en paralelo: según el autor, para expresar la acción paralela el
español, como gran parte de las lenguas indoeuropeas, gramaticalizó la forma
\forma{tanto X como Y}, heredera de una estructura latina que también originó
esquemas comparativos. Señala también que esta forma gramaticalizada tiene un
correlato negativo: \forma{ni X ni Y}.} \forma{ni\ldots\ ni\ldots;
tanto\ldots\ cuanto\ldots; así\ldots\ como}. Cuando se trata de una
coordinación múltiple, como ya hemos anticipado, la conjunción \forma{y}
normalmente se coloca solo ante el último elemento, pero se interpreta un
vínculo de suma entre todos los elementos del conjunto.

\begin{destacado}
\textbf{Para tener en cuenta.} La omisión de la conjunción ante el último
elemento o bien su repetición ante cada elemento coordinado genera efectos que
dan lugar a dos figuras retóricas. Como se explica en la \textit{NGLE}
\parencite[§\,31.2b]{rae2025}, el proceso de coordinar expresiones sin
conjunción manifiesta se denomina \termino{asíndeton} y constituye una de las
formas de yuxtaposición; la suspensión de la conjunción final deja a menudo en
suspenso la enumeración. En este caso, se interpreta que la enumeración no es
exhaustiva: \forma{Compramos cuadernos, hojas, lápices, lapiceras\ldots} En
cambio, se denomina \termino{polisíndeton} la figura retórica que consiste en el
proceso contrario, es decir, la repetición de la conjunción copulativa ante cada
miembro coordinado: se trata de un recurso enfático dirigido a lograr mayor
energía en la expresión, como en \forma{Se lo dije una, y dos, y tres veces, pero
no me hacía caso} \parencite[§\,31.2c]{rae2025}.
\end{destacado}

\paragraph{Lectura distributiva y colectiva de la coordinación copulativa.}\label{sec:12-distributiva}
En general, cuando se construyen predicados en los que participan sintagmas
nominales unidos por una conjunción copulativa, emerge una \termino{lectura
distributiva}, de modo que a cada coordinando le corresponde la misma
predicación, por lo que \exref{ex:12-teatro-a} puede parafrasearse como
\exref{ex:12-teatro-b}.

\begin{exe}
  \ex
    \begin{xlist}
      \ex\label{ex:12-teatro-a} María \textit{y} Ana fueron al teatro.
      \ex\label{ex:12-teatro-b} María fue al teatro \textit{y} Ana fue al
          teatro.
    \end{xlist}
\end{exe}

Por otro lado, cuando tenemos predicados de naturaleza colectiva (véase el
recuadro \textit{Para ampliar}), como en \exref{ex:12-parecen} y
\exref{ex:12-cunadas}, la presencia de una coordinación requiere de una
interpretación conjunta o acumulativa (\termino{lectura colectiva}), y se rechaza
la lectura distributiva.

\begin{exe}
  \ex\label{ex:12-parecen}
    \begin{xlist}
      \ex Pedro y Joaquín se parecen.
      \ex[*]{\label{ex:12-parecen-b} Pedro se parece \textit{y} Joaquín se
          parece.}
    \end{xlist}
  \ex\label{ex:12-cunadas}
    \begin{xlist}
      \ex María \textit{y} Ana son cuñadas.
      \ex[*]{\label{ex:12-cunadas-b} María es cuñada \textit{y} Ana es cuñada.}
    \end{xlist}
\end{exe}

\begin{recuadro}[Para ampliar]
De acuerdo con la \textit{NGLE} \parencite[§\,31.6c]{rae2025}, los predicados de
naturaleza colectiva «se suelen llamar simétricos y se caracterizan por
predicarse de grupos, manifestados mediante plurales o elementos coordinados».
Estos predicados pueden ser verbos (\forma{casarse, conversar, juntar,
negociar}), adjetivos (\forma{coincidente, compatible, diferente}) o sustantivos
(\forma{adversario, amigo, cuñado, pariente, socio}). También tienen
características de predicados simétricos un SP con \forma{entre} como núcleo y
algunos adverbios o locuciones adverbiales, como \forma{lejos, cerca,
conjuntamente, al lado, en colaboración}.
\end{recuadro}

Como resulta claro a partir de la agramaticalidad de \exref{ex:12-parecen-b} y
\exref{ex:12-cunadas-b}, el predicado no puede repetirse para cada elemento del
SConj, por lo que la única lectura posible es la colectiva: se predica de los
coordinandos en forma conjunta. Tampoco hay ambigüedad cuando se emplean, en
cualquier predicado, conjunciones copulativas compuestas, también llamadas
\textit{discontinuas} o \textit{correlativas}, porque impulsan una
interpretación distributiva. En este sentido, la agramaticalidad de una oración
como \exref{ex:12-tanto} se debe a la tensión que se produce entre la lectura
colectiva del predicado \forma{ser cuñadas} y la lectura distributiva que
requiere la conjunción \forma{tanto\ldots\ como}:

\begin{exe}
  \ex[*]{\label{ex:12-tanto} Tanto María como Ana son cuñadas entre sí.}
\end{exe}

Sin embargo, en predicados no simétricos, es posible encontrar ejemplos de
ambigüedad entre la interpretación colectiva y la distributiva de \forma{y}. Por
ejemplo, una oración como \exref{ex:12-pastel-a} puede interpretarse como
\exref{ex:12-pastel-b} o como \exref{ex:12-pastel-c}:

\begin{exe}
  \ex
    \begin{xlist}
      \ex\label{ex:12-pastel-a} María \textit{y} Ana cocinaron un pastel de papas
          y un budín.
      \ex\label{ex:12-pastel-b} María \textit{y} Ana cocinaron un pastel de papas
          y un budín \textit{juntas}. (Lectura colectiva).
      \ex\label{ex:12-pastel-c} María cocinó un pastel de papas y un budín,
          \textit{y} Ana cocinó un pastel de papas y un budín. (Lectura
          distributiva).
      \ex\label{ex:12-pastel-d} María \textit{y} Ana cocinaron un pastel de papas
          y un budín, \textit{respectivamente}. (Lectura ordenada).
    \end{xlist}
\end{exe}

Asimismo, como vemos en \exref{ex:12-pastel-d}, la ambigüedad puede evitarse
mediante la incorporación de \forma{respectivamente}.

\begin{ejercicio}\label{ej:12-lecturas}
Lea las siguientes oraciones y proponga las lecturas posibles, según se muestra
en el ejemplo resuelto~1.

\begin{exe}
  \ex Frodo y Sam compraron frutas y galletas para el viaje.
  \ex Frodo y Sam se llevan bien desde la infancia.
  \ex Frodo y Sam viajan a Mordor.
  \ex Sam y Gollum se odian.
  \ex Después de la batalla, Frodo y Aragorn viajan hacia Mordor y hacia el
      Bosque de Fangorn.
\end{exe}

\begin{resuelto}
1. a) Lectura distributiva: \forma{Frodo compró frutas y galletas para el viaje
y Sam también compró frutas y galletas para el viaje} / \forma{Tanto Frodo como
Sam compraron frutas y galletas}. b) Lectura colectiva: \forma{Frodo y Sam
compraron frutas y galletas juntos}. c) Lectura ordenada: \forma{Frodo y Sam
compraron frutas y galletas, respectivamente}.
\end{resuelto}
\end{ejercicio}

\paragraph{Coordinación simétrica y asimétrica.}\label{sec:12-simetrica}
Por último, retomemos un aspecto de los coordinandos que ya habíamos mencionado
antes, por lo que repetiremos los ejemplos de \exref{ex:12-hermano4},
\exref{ex:12-maria}, \exref{ex:12-cansado} y \exref{ex:12-bebida}, reformulados
como \exref{ex:12-sim-a}--\exref{ex:12-sim-d}:

\begin{exe}
  \ex
    \begin{xlist}
      \ex\label{ex:12-sim-a} Mi hermano lavó los platos y estaba contento.
      \ex María llegó a la oficina y encendió la computadora.
      \ex Entró cansado y sin prisa.
      \ex\label{ex:12-sim-d} Compramos la bebida y todo lo que estaba en oferta.
    \end{xlist}
\end{exe}

En principio, la coordinación copulativa establece una relación simétrica entre
los coordinandos, que se manifiesta no solo en la homogeneidad categorial y
funcional, que hemos visto en §\ref{sec:12-coordinacion}, sino también en el
orden: si la relación es simétrica, los coordinandos son intercambiables. Ahora
bien, como afirma \textcite[282]{ditullio2014}, cuando tal intercambio no es
posible, «es porque se asocian al coordinante otros valores semánticos» y se
produce lo que se conoce como \termino{coordinación asimétrica}, en la que los
coordinandos no son intercambiables. Veamos en \exref{ex:12-inv} los resultados
del cambio de orden de las oraciones anteriores.

\begin{exe}
  \ex\label{ex:12-inv}
    \begin{xlist}
      \ex\label{ex:12-inv-a} Mi hermano estaba contento \textit{y} lavó los
          platos.
      \ex[*]{\label{ex:12-inv-b} María encendió la computadora \textit{y} llegó
          a la oficina.}
      \ex\label{ex:12-inv-c} Entró sin prisa \textit{y} cansado.
      \ex\label{ex:12-inv-d} Compramos todo lo que estaba en oferta \textit{y} la
          bebida.
    \end{xlist}
\end{exe}

Notamos que, luego del cambio de orden, la gramaticalidad se mantiene
indiscutiblemente en \exref{ex:12-inv-a}, \exref{ex:12-inv-c} y
\exref{ex:12-inv-d}, por lo que validamos que se trata de coordinaciones
simétricas. Tal vez podríamos señalar que, con la inversión,
\exref{ex:12-inv-a} adquiere cierto matiz de causa, en tanto podría entenderse
como equivalente a \forma{Como mi hermano estaba contento, lavó los platos}. En
cambio, \exref{ex:12-inv-b} es agramatical,\pendiente{Con la lectura de
sucesión, \exref{ex:12-inv-b} es gramatical, aunque describe otra secuencia de
eventos; quizá corresponda \# en lugar de *} lo que nos lleva a afirmar que los
coordinandos no son simétricos en su configuración, es decir, no se puede
alterar el orden de aparición. En este caso, notamos que el orden de los
constituyentes se encuentra asociado a una relación temporal que no permite la
inversión: María primero llega a la oficina y luego enciende la computadora para
trabajar.

Además del valor temporal, la conjunción copulativa puede adquirir otros valores
semánticos: condicional, concesivo, adversativo, ilativo. Tal como se explica en
la \textit{NGLE} \parencite[§\,31.2m]{rae2025}, cuando uno de los dos predicados
coordinados no expresa un suceso, sino una situación o un estado de cosas, no se
obtiene la interpretación temporal. En esos casos son comunes la interpretación
adversativa (\forma{Tenía examen y no era capaz de recordar nada}, es decir,
\sig{y sin embargo\ldots}), la ilativa (\forma{Estaba cansado y me fui a
dormir}, es decir, \sig{luego\ldots}) o la condicional, frecuente en los
imperativos retóricos (\forma{Cásate y verás}).

Respecto de estos valores asociados a la conjunción, algunos autores consideran
que se trata, en realidad, de un fenómeno vinculado con aspectos pragmáticos, es
decir, con el uso de la lengua. Por ejemplo, \textcite[195]{jimenezjulia2023}
afirma que la posibilidad de interpretar contextualmente otros valores «no es una
consecuencia del valor de \textit{y}, sino de nuestro conocimiento de la
realidad» y de la relación lógica esperable cuando se encadenan, en cierto
orden, los eventos expresados.

\begin{ejercicio}\label{ej:12-copulativas}
Lea las siguientes oraciones, adaptadas del artículo «\textit{Homo sapiens}, el
animal que todo lo extingue»,\footnote{P. Esteban, «\textit{Homo sapiens}, el
animal que todo lo extingue», \textit{Página/12}, 12/10/2024
(\url{https://www.pagina12.com.ar/774085-homo-sapiens-el-animal-que-todo-lo-extingue}).}
y luego resuelva lo pedido, según se muestra en el ejemplo resuelto~1.

\begin{consignas}
  \item Encierre entre corchetes las coordinaciones copulativas e indique, en cada
        caso, la función sintáctica del SConj y la naturaleza categorial de los
        elementos coordinados.
  \item Distinga entre coordinación simétrica y asimétrica, cambiando el orden de
        los coordinandos. En caso de que haya asimetría, identifique el valor
        semántico asociado a la conjunción.
\end{consignas}

\begin{exe}
  \ex Las especies exóticas introducidas por los humanos suelen ser
      determinantes para su eliminación porque, como están en islas, no pueden
      escapar ni pueden asentarse con éxito en otras latitudes.
  \ex Meteoritos y volcanes, en el pasado, provocaron las extinciones masivas
      documentadas a lo largo de la historia.
  \ex El dodo vivió en la Isla Mauricio y se extinguió hacia el siglo XVII.
  \ex La compañía aislará y cultivará óvulos y células productoras de esperma
      para luego editar las secuencias de ADN que coincidan con las del dodo.
  \ex Estas células darían lugar a un animal mezcla de dodo y paloma.
  \ex El mamut lanudo habitó Asia y Europa hace 6 mil años, y podría volver a
      existir gracias a la ciencia.
\end{exe}

\begin{resuelto}
1. a) Las especies exóticas introducidas por los humanos suelen ser
determinantes para su eliminación porque, como están en islas,
[\etq{SConj/Pred} [\etq{SV/Esp} no pueden escapar] ni\etq{Conj/n}
[\etq{SV/Comp} pueden asentarse con éxito en otras latitudes]]. b) La
coordinación es asimétrica. No es posible cambiar el orden de los coordinandos,
debido a la presencia de un valor temporal de sucesión asociado a la
conjunción.\pendiente{La inversión (\textit{no pueden asentarse\ldots\ ni pueden
escapar}) parece aceptable: revisar el carácter asimétrico}
\end{resuelto}
\end{ejercicio}

\subsubsection{La coordinación disyuntiva}\label{sec:12-disyuntiva}

La coordinación disyuntiva se expresa, principalmente, a partir del uso de la
conjunción simple \forma{o}, que puede unir dos o más elementos que se
interpretan como \sig{opciones}. Además, se pueden utilizar conjunciones
discontinuas: \forma{o\ldots\ o\ldots, bien\ldots\ bien\ldots, ya\ldots\
ya\ldots, ora\ldots\ ora\ldots, sea\ldots\ sea\ldots, fuera\ldots\ fuera\ldots}
Excepto \forma{o\ldots\ o}, las restantes, debido a su alto grado de formalidad,
en la actualidad tienen uso restringido y se emplean sobre todo en la lengua
escrita.

Tomando, entonces, como ejemplo típico el de la conjunción \forma{o}, podemos
señalar que se establece un primer matiz en la interpretación de la disyunción en
función de si la elección entre las opciones es obligatoria o no. En el primer
caso, hablamos de una \termino{interpretación exclusiva}, como en
\exref{ex:12-disy-a} y \exref{ex:12-disy-b}. En cambio, cuando la elección no se
impone necesariamente, hablamos de una \termino{interpretación inclusiva}, en la
que las opciones presentadas se interpretan como instancias posibles, como en
\exref{ex:12-disy-c} y \exref{ex:12-disy-d}:

\begin{exe}
  \ex
    \begin{xlist}
      \ex\label{ex:12-disy-a} La azafata preguntó a los pasajeros si querían
          carne \textit{o} pasta.
      \ex\label{ex:12-disy-b} ¿Compramos el pasaje ahora \textit{o} esperamos a
          que baje el precio?
      \ex\label{ex:12-disy-c} Podemos ir al cine mañana, el sábado \textit{o} el
          domingo.
      \ex\label{ex:12-disy-d} Mañana podemos ir al cine, \textit{o} visitar unos
          amigos \textit{o} preparar una rica cena.
    \end{xlist}
\end{exe}

\begin{recuadro}[Para ampliar]
Con respecto a la interpretación inclusiva y la exclusiva de la conjunción
\forma{o} del español, \textcite[195]{jimenezjulia2023} establece un vínculo con
los valores expresados en latín por las conjunciones \forma{aut} y \forma{uel}.
De acuerdo con lo señalado por este autor, en latín \forma{aut} expresaba el
valor exclusivo y \forma{uel} (o \forma{siue / seu}), el inclusivo. El español
heredó la conjunción \forma{o} de la primera, pero no conservó \forma{uel}. Ello,
sin embargo, no implicó la pérdida de los dos valores: se produjo una
reconversión, de modo que la conjunción \forma{o} pasó a ser la general para
ambos valores. Asimismo, Jiménez Juliá añade que la repetición de la conjunción,
que la convierte en correlativa (\forma{o X o Y}), solo expresa valor exclusivo.
\end{recuadro}

En algunos casos, la conjunción \forma{o} tiene otras interpretaciones. Por un
lado, puede usarse con un sentido denominativo para introducir una aclaración
del valor de una expresión. Como vemos en los ejemplos de
\exref{ex:12-carpincho}, en estas instancias se trata de un uso metalingüístico:

\begin{exe}
  \ex\label{ex:12-carpincho}
    \begin{xlist}
      \ex El carpincho \textit{o} capibara es el roedor más grande del mundo.
      \ex Pekín \textit{o} Beijing es la capital de la República Popular China.
    \end{xlist}
\end{exe}

Por otro lado, la disyunción puede usarse con un valor correctivo, con el
objetivo de introducir una reformulación que, a su vez, incorpora algún matiz
respecto del primer coordinando:

\begin{exe}
  \ex
    \begin{xlist}
      \ex Nueva Constitución, \textit{o} el cuento de la inversión extranjera.
          (\textit{Página/12}, 2/9/2024)
      \ex Tu obra fue muy corta \textit{o}, mejor dicho, el tiempo se pasó
          volando.
    \end{xlist}
\end{exe}

Además, la conjunción \forma{o} puede usarse de forma discontinua o
distributiva, es decir, antepuesta a ambos elementos de la coordinación:

\begin{exe}
  \ex
    \begin{xlist}
      \ex Podemos ir al concierto \textit{o} el jueves \textit{o} el viernes
          (pero no los dos días).
      \ex \textit{O} nos unimos al mundo globalizado \textit{o} quedamos
          aislados.
    \end{xlist}
\end{exe}

Si bien hay otras conjunciones disyuntivas discontinuas, como \forma{ya\ldots\
ya\ldots, bien\ldots\ bien\ldots, ora\ldots\ ora\ldots}, su frecuencia de uso es
mucho menor. Respecto de estas conjunciones correlativas,
\textcite[196]{jimenezjulia2023} señala que no son un tipo específico de
conjunciones, sino que se les da ese nombre a estructuras coordinadas
copulativas o disyuntivas expresadas con nexos correlativos. Como señala este
autor, las conjunciones correlativas pueden tener un valor disyuntivo, como en
\exref{ex:12-corr-a}, o copulativo, como en \exref{ex:12-corr-b}:

\begin{exe}
  \ex
    \begin{xlist}
      \ex\label{ex:12-corr-a} Todos sabemos que \textit{prefijo} es lo que
          \sig{figura o se coloca delante} de un número o una palabra, ya sea el
          \textit{34} para llamar a España desde el extranjero, ya sea el
          elemento compositivo \textit{auto-}. (Fundéu RAE)
      \ex\label{ex:12-corr-b} La perfección absoluta, bien en las artes, bien en
          las demás cosas, no es de aquellas que el cielo quiso conceder a la
          tierra. (M.~Menéndez Pelayo, \textit{Historia})
    \end{xlist}
\end{exe}

Como puede observarse en \exref{ex:12-corr-b}, cuando el valor es copulativo,
las conjunciones correlativas son equivalentes a la expresión
\forma{tanto\ldots\ como\ldots}. En caso de que su valor sea disyuntivo, solo se
establece la interpretación de una disyunción exclusiva, es decir, se impone
necesariamente una decisión por uno de los dos elementos coordinados.

\begin{ejercicio}\label{ej:12-disyuntivas}
Lea las siguientes oraciones y luego resuelva lo pedido según se muestra en el
ejemplo resuelto~1.

\begin{consignas}
  \item Encierre entre corchetes las coordinaciones disyuntivas y determine la
        función sintáctica del SConj y la naturaleza categorial de los
        coordinandos.
  \item Identifique el valor con el que se interpreta la coordinación.
\end{consignas}

\begin{exe}
  \ex El carpincho o capibara es el roedor más grande del mundo.
  \ex Los capibaras son dóciles y usualmente los humanos los acarician o les dan
      de comer.
  \ex Dependiendo de la subespecie, el periodo de gestación dura unos 110 días o
      unos 150 días.
  \ex Sus depredadores naturales más importantes son félidos, como jaguares o
      pumas.
  \ex En la lengua tupí, el término es \textit{ka'apiûara}, que significa \sig{el
      que come hojas delgadas}, o \sig{comedor de hierba}.
  \ex En 2020, los carpinchos invadieron Nordelta o, mejor dicho,\footnote{Como
      se verá en §\ref{sec:12-marcadores}, \forma{mejor dicho} se considera un
      conector discursivo, que se combina con la conjunción, pero no es parte de
      ella.} volvieron a su territorio autóctono en los humedales del río
      Paraná.
\end{exe}

\begin{resuelto}
1. a) [\etq{SConj/Suj} [\etq{SN/Esp} El carpincho] o\etq{Conj/n}
[\etq{SN/Comp} capibara]] es el roedor más grande del mundo. b) Disyunción
metalingüística denominativa: introduce una segunda denominación del término
\forma{carpincho}.
\end{resuelto}
\end{ejercicio}

\subsubsection{La coordinación adversativa}\label{sec:12-adversativa}

En la coordinación adversativa se establece entre los coordinandos una relación
de \sig{oposición}. Las conjunciones adversativas principales son \forma{pero} y
\forma{sino}. También es adversativa la conjunción
\forma{mas},\footnote{\textcite[197]{jimenezjulia2023} comenta que \forma{mas}
proviene del latín \forma{magis}, origen de la conjunción adversativa dominante
en lenguas románicas como el francés (\forma{mais}) o el portugués
(\forma{mas}); en cambio, tanto en el español europeo como en las variedades
americanas parece relegada a registros cultos.} de características similares a
\forma{pero}, aunque su uso está más restringido por motivos de registro. Dado
que las dos conjunciones principales tienen características formales bien
diferenciadas, las presentaremos por separado.

\paragraph{La conjunción \textit{pero}.}\label{sec:12-pero}
Se puede utilizar para coordinar sintagmas de distinto tipo, aunque es más
frecuentemente utilizada para coordinar oraciones:

\begin{exe}
  \ex
    \begin{xlist}
      \ex\label{ex:12-depto-a} El departamento es [\etq{SA} chico], \textit{pero}
          [\etq{SA} luminoso].
      \ex El departamento tiene [\etq{SN} muchas puertas], \textit{pero}
          [\etq{SN} pocas ventanas].
      \ex\label{ex:12-depto-c} {}[\etq{SF/O} El departamento está a buen precio],
          \textit{pero} [\etq{SF/O} aún no apareció ningún comprador].
    \end{xlist}
\end{exe}

Con respecto a los valores de la relación adversativa establecida con la
conjunción \forma{pero}, \textcite[284]{ditullio2014} señala la existencia, por
un lado, de un valor de \termino{contraexpectativa} y, por el otro, de una
\termino{oposición valorizadora}. En el primer caso, se establece una relación de
contraposición entre lo expresado por ambos coordinandos, pero a partir de la
interpretación de una inferencia del hablante. Por ejemplo, en
\exref{ex:12-depto-c}, de lo dicho por el primer coordinando (\forma{el
departamento está a buen precio}) se podría inferir que se debería vender
rápido, pero esa inferencia se contrapone con lo manifestado por el segundo
coordinando (\forma{aún no apareció ningún comprador}), en el que se afirma un
evento o un estado de cosas contrario al inferido. En palabras de Di Tullio, la
conjunción \forma{pero} cancela la expectativa generada por el primer
coordinando.

Con respecto a la lectura valorizadora, la autora señala que tal interpretación
se produce cuando los elementos coordinados se pueden relacionar con una
jerarquía, de modo que la oposición permite priorizar lo expresado en el segundo
coordinando. Por ejemplo, en \exref{ex:12-depto-a} se contraponen dos
características de un departamento: \forma{ser chico} y \forma{ser luminoso}.
Pero, por el orden dado a los elementos, se interpreta una valorización de la
segunda característica y, por ende, una visión positiva del departamento. Si
alteramos el orden de los coordinandos (\forma{El departamento es luminoso, pero
chico}), la característica ahora valorizada es \forma{ser chico}, la cual suele
asociarse con un rasgo negativo, por lo que la oración se interpreta como un
juicio negativo sobre la propiedad.

En ambos casos, se destaca el hecho de que la coordinación adversativa se vincula
necesariamente con las inferencias que el hablante puede establecer a partir
tanto de su conocimiento del mundo como del contexto discursivo. Es en este
sentido que, como se afirma en la \textit{NGLE} \parencite[§\,31.10i]{rae2025},
la coordinación adversativa obliga al oyente a establecer relaciones de causa y
efecto entre oraciones mediante mecanismos discursivos que no se basan
únicamente en el contenido de las palabras. A modo de ejemplo, veamos cómo
funciona la coordinación adversativa en las siguientes oraciones de \textit{Éste
es el mar}, de Mariana Enriquez:

\begin{exe}
  \ex
    \begin{xlist}
      \ex\label{ex:12-helena} Helena las miraba atenta y tímida, \textit{pero}
          finalmente se atrevió a hablar.
      \ex\label{ex:12-humanos} Cuando los humanos investigaban los Misterios de
          los Últimos Días, notaban la rareza, \textit{pero} no comprendían su
          enormidad.
      \ex\label{ex:12-enjambre} El Enjambre era vasto, eran muchas sus
          integrantes, \textit{pero} no podían cambiar de rostro ni de cuerpo.
    \end{xlist}
\end{exe}

En \exref{ex:12-helena} se muestra el valor de contraexpectativa, dado que el
segundo coordinando (\forma{finalmente se atrevió a hablar}) cancela la
expectativa que genera el primero, a partir de la inferencia de que la timidez
no iba a dejar que hable. De manera similar, en \exref{ex:12-humanos}, a partir
del primer coordinando se puede inferir que los humanos, dado que pueden notar la
rareza, podrían encontrarle una explicación; sin embargo, tal expectativa se ve
cancelada por el contenido del segundo coordinando. Finalmente, en
\exref{ex:12-enjambre}, predomina una interpretación valorizadora, dado que los
coordinandos expresan dos características de las integrantes del Enjambre: ser
numerosas y no poder cambiar de rostro ni de cuerpo. Dado el orden elegido, en
este fragmento la evaluación se interpreta como negativa, en tanto se realza
aquello que las integrantes no pueden hacer. Dicha valoración podría volverse
positiva con simplemente alterar el orden de los coordinandos: \forma{No podían
cambiar de rostro ni de cuerpo, pero el Enjambre era vasto, eran muchas sus
integrantes}.

\paragraph{La conjunción \textit{sino}.}\label{sec:12-sino}
Esta conjunción, al igual que \forma{pero}, establece una coordinación binaria
opositiva, con la particularidad de que exige la presencia de una negación en el
primer coordinando y el segundo coordinando se interpreta como el reemplazo del
elemento negado, por lo que la coordinación constituye una \termino{oposición
exclusiva}. La conjunción \forma{sino} puede coordinar distintos tipos de
sintagmas, desde sintagmas léxicos hasta oraciones independientes o
subordinadas:

\begin{exe}
  \ex\label{ex:12-sino}
    \begin{xlist}
      \ex Juan llegó \textit{no} a la mañana, \textit{sino} a la tarde.
      \ex Pedro \textit{no} tiene dos perros, \textit{sino} dos gatos.
      \ex El animal que trajiste \textit{no} es macho, \textit{sino} hembra.
      \ex Ana \textit{no} dijo que compró el libro, \textit{sino} que lo había
          visto a la venta.
    \end{xlist}
\end{exe}

Cabe destacar que, en muchas ocasiones, la negación del primer término exigida
por \forma{sino} se usa de forma combinada con el adverbio \forma{solo},
correlativo del adverbio \forma{también}, lo cual cambia la interpretación de
oposición exclusiva, de modo que el segundo término de la coordinación ya no se
interpreta como un reemplazo, sino que adquiere valor de adición. Para que tal
interpretación sea posible, es fundamental que haya compatibilidad entre los
predicados de cada coordinando. Como vemos, si agregamos \forma{solo} y
\forma{también} a los ejemplos de \exref{ex:12-sino}, podemos obtener secuencias
gramaticales, como \exref{ex:12-solo-b}, o dudosas, como \exref{ex:12-solo-d},
ambas con cambio de significado: en \exref{ex:12-solo-b}, Pedro pasa a tener
cuatro mascotas en vez de dos y, en \exref{ex:12-solo-d}, se pierde el sentido
exclusivo de la adversativa. Por último, en \exref{ex:12-solo-a} y
\exref{ex:12-solo-c} las secuencias que se obtienen son incoherentes, dado que
los significados de los elementos léxicos coordinados ---\forma{mañana / tarde} y
\forma{macho / hembra}--- son incompatibles y no pueden añadirse.

\begin{exe}
  \ex
    \begin{xlist}
      \ex[\#]{\label{ex:12-solo-a} Juan llegó \textit{no solo} a la mañana,
          \textit{sino también} a la tarde.}
      \ex\label{ex:12-solo-b} Pedro \textit{no solo} tiene dos perros,
          \textit{sino también} dos gatos.
      \ex[\#]{\label{ex:12-solo-c} El animal que trajiste \textit{no solo} es
          macho, \textit{sino también} hembra.}
      \ex[?]{\label{ex:12-solo-d} Ana \textit{no solo} dijo que compró el libro,
          \textit{sino también} que lo había visto a la venta.}
    \end{xlist}
\end{exe}

\begin{ejercicio}\label{ej:12-adversativas}
Lea las siguientes oraciones, adaptadas del artículo «\textit{Homo sapiens}, el
animal que todo lo extingue», y luego resuelva lo pedido según se muestra en el
ejemplo resuelto~1.

\begin{consignas}
  \item Encierre entre corchetes las coordinaciones adversativas y determine la
        función sintáctica del SConj y la naturaleza categorial de los
        coordinandos.
  \item Analice el valor con el que se interpreta la coordinación.
\end{consignas}

\begin{exe}
  \ex Las sociedades no solo barren con el número de especies, sino que arrancan
      de raíz ecosistemas enteros.
  \ex Estos números representan una pérdida ya consolidada del 5 por ciento de
      los linajes, es decir, de especies distintas, pero con un ancestro en
      común.
  \ex No solo se achicará la biodiversidad de la fauna, sino que también se
      perderán funciones ecosistémicas cruciales, como la polinización y el
      control de plagas.
  \ex La paloma migratoria era la especie de ave más abundante en Norteamérica,
      pero ahora está completamente extinta.
  \ex El ser humano no se siente parte del ecosistema, sino que está totalmente
      desconectado de la naturaleza de la que forma parte.
  \ex Hay medidas para frenar este conflicto, pero se tienen que tomar muy
      rápido.
\end{exe}

\begin{resuelto}
1. a) [\etq{SConj/OC} [\etq{SF/Esp} Las sociedades no solo barren con el número
de especies], sino que\etq{Conj/n} [\etq{SF/Comp} arrancan de raíz ecosistemas
enteros]]. b) La coordinación adversativa con \forma{sino} no se interpreta como
oposición exclusiva, sino como adición contrastiva, por la presencia del
adverbio \forma{solo} en el primer coordinando.
\end{resuelto}
\end{ejercicio}

\subsection{Diferencias entre conjunciones y marcadores discursivos}\label{sec:12-marcadores}

Las conjunciones coordinantes a menudo se utilizan combinadas con elementos como
\forma{además, sin embargo, encima}. Si bien estos elementos hacen su aporte para
establecer la conexión entre los constituyentes, no pertenecen a la clase de las
conjunciones coordinantes, sino que se los conoce como \termino{marcadores} o
\termino{conectores discursivos}.

En el \textit{Glosario de términos gramaticales}, se define el marcador
discursivo como una «unidad léxica, formada por una o varias palabras, que
establece vínculos lógicos, discursivos o argumentativos entre los segmentos de
un texto» \parencite[74]{bosque2019}. La mayor parte de los marcadores
discursivos son locuciones adverbiales (\forma{sin embargo, por si fuera poco,
no obstante}), pero dentro del paradigma también encontramos adverbios
(\forma{encima, además}), interjecciones (\forma{claro, bueno}) y locuciones
conjuntivas (\forma{ya que, puesto que}). Como se explica en la \textit{NGLE}
\parencite[§\,30.12a]{rae2025}, el concepto de conector o marcador discursivo
«no constituye una clase sintáctica de palabras análoga a \textit{verbo,
conjunción} o \textit{adverbio}, sino un grupo establecido con criterios
textuales».

Para dar cuenta de la necesidad de diferenciar las conjunciones de los conectores
discursivos, debemos tener en cuenta que las conjunciones tienen una posición fija
entre los coordinandos, como se observa en \exref{ex:12-fija}, mientras que los
conectores pueden desplazarse, tal como se muestra en \exref{ex:12-movil}:

\begin{exe}
  \ex\label{ex:12-fija}
    \begin{xlist}
      \ex Juan \textit{y} Pedro vendrán mañana.
      \ex[*]{Juan Pedro \textit{y} vendrán mañana.}
    \end{xlist}
  \ex\label{ex:12-movil}
    \begin{xlist}
      \ex Juan salió temprano y, \textit{sin embargo}, no llegó a tiempo.
      \ex Juan salió temprano y no llegó a tiempo, \textit{sin embargo}.
    \end{xlist}
\end{exe}

Una segunda diferencia importante entre marcadores y conjunciones radica en que,
cuando están al inicio de la oración, los primeros van seguidos de pausa,
mientras que las conjunciones generalmente forman una unidad entonacional con el
constituyente al que preceden. En las oraciones de \exref{ex:12-pausa} puede
apreciarse el contraste:

\begin{exe}
  \ex\label{ex:12-pausa}
    \begin{xlist}
      \ex \textit{Además}, Juan todavía no llegó.
      \ex[*]{\textit{Pero}, Juan todavía no llegó.}
      \ex \textit{Pero} Juan todavía no llegó.
    \end{xlist}
\end{exe}

Otro aspecto que muestra la necesidad de diferenciar entre conjunciones
coordinantes y marcadores discursivos es la posibilidad de combinación: al
asociar dos conjunciones coordinantes se obtiene una secuencia agramatical, como
se muestra en \exref{ex:12-comb-a}, mientras que la combinación de una conjunción
y un conector discursivo produce secuencias que sí son gramaticales, como la de
\exref{ex:12-comb-b}:

\begin{exe}
  \ex
    \begin{xlist}
      \ex[*]{\label{ex:12-comb-a} Juan todavía no llegó, \textit{pero y} podemos
          empezar.}
      \ex\label{ex:12-comb-b} Juan todavía no llegó, \textit{pero, sin embargo},
          podemos empezar.
    \end{xlist}
\end{exe}

En definitiva, los marcadores discursivos sirven para orientar «la manera en que
la oración o el fragmento oracional sobre el que inciden han de ser
interpretados» en relación con el contexto precedente o con las inferencias que
de él se obtienen \parencite[§\,30.12i]{rae2025}, pero no tienen las mismas
propiedades sintácticas que las conjunciones. De hecho, un contraste adicional
puede encontrarse en el inventario de ítems de cada grupo. Como ya vimos, las
conjunciones coordinantes, si bien pueden presentar distintos valores
(copulativo, disyuntivo, adversativo), constituyen un inventario bastante
acotado: \forma{y, ni, o, pero, sino} (y las variantes correlativas). En cambio,
existe un amplio repertorio de marcadores discursivos, tan abundante que
normalmente se presenta en subconjuntos, agrupados por su función discursiva. En
la tabla~\ref{tab:12-marcadores} presentamos la propuesta de clasificación de
los grupos fundamentales, tomada de la \textit{NGLE}
\parencite[§\,30.13a]{rae2025}.

{\small
\begin{longtable}{@{}>{\raggedright\arraybackslash}p{38mm}
    >{\raggedright\arraybackslash\itshape}p{85mm}@{}}
  \caption{Marcadores discursivos adverbiales}\label{tab:12-marcadores}\\
  \toprule
  Grupo & \normalfont Marcadores \\
  \midrule
  \endfirsthead
  \toprule
  Grupo & \normalfont Marcadores \\
  \midrule
  \endhead
  1. Aditivos y de precisión o particularizadores & a decir verdad, además,
    análogamente, aparte, asimismo, de hecho, encima, en el fondo, en realidad,
    es más, por añadidura, por otro lado, por si fuera poco, sobre todo \\
  2. Adversativos y contraargumentativos & ahora bien, al contrario, antes bien,
    después de todo, empero, en cambio, eso sí, no obstante, por el contrario,
    sin embargo, todo lo contrario \\
  3. Concesivos & así y todo, aun así, con todo, de cualquier manera, de todas
    formas, de todas maneras, de todos modos, en cualquier caso \\
  4. Consecutivos e ilativos & así pues, consiguientemente, de ese modo, en
    consecuencia, entonces, por consiguiente, por ende, por lo tanto, por tanto,
    pues \\
  5. Explicativos & a saber, es decir, esto es, o sea \\
  6. Reformuladores & dicho con otras palabras, dicho en otros términos, dicho de
    otra forma, dicho de otra manera, de otro modo, más claramente, más
    llanamente, hablando claro \\
  7. Ejemplificativos & así, así por ejemplo, así tenemos, por ejemplo,
    verbigracia \\
  8. Rectificativos & más bien, mejor dicho, por mejor decir \\
  9. Recapitulativos & a fin de cuentas, al fin y al cabo, en conclusión, en
    definitiva, en fin, en resumen, en resumidas cuentas, en síntesis, en suma,
    en una palabra, resumiendo, total \\
  10. Ordenadores & a continuación, antes que nada, de una parte, de otra parte,
    en primer lugar, en primer término, finalmente, para empezar, para
    terminar \\
  11. De apoyo argumentativo & así las cosas, dicho esto, en vista de ello, pues
    bien \\
  12. De digresión & a propósito, a todo esto, dicho sea de paso, entre
    paréntesis, por cierto \\
  \bottomrule
\end{longtable}}

\subsection{Coordinación y elipsis}\label{sec:12-elipsis}

Un fenómeno estrechamente vinculado con la coordinación de sintagmas y oraciones
es la \termino{elipsis}, es decir, la posibilidad de eliminar elementos del
segundo coordinando, en función de que pueden interpretarse a partir del primer
coordinando. En algunos casos, la elipsis es fácilmente detectable: cuando se
omite el verbo del segundo coordinando, como en \exref{ex:12-elip-a}, donde se
puede observar la presencia de una coma que marca la elisión; cuando se elide el
verbo y algún complemento, como en \exref{ex:12-elip-b}, o cuando en el segundo
coordinando se emplean los adverbios \forma{también, tampoco, sí} o \forma{no},
que dan lugar a la elipsis de todo el sintagma verbal, como en
\exref{ex:12-elip-c} y \exref{ex:12-elip-d}:\pendiente{En
\exref{ex:12-elip-b} se elide el verbo con su COD, pero el remanente es un
adjunto}

\begin{exe}
  \ex\label{ex:12-elip}
    \begin{xlist}
      \ex\label{ex:12-elip-a} Juan estudia Letras \textit{y} Pedro, Filosofía.
      \ex\label{ex:12-elip-b} Ana leyó el libro en su casa \textit{y} María, en la
          plaza.
      \ex\label{ex:12-elip-c} Juan no estudia Filosofía, \textit{pero} Pedro sí.
      \ex\label{ex:12-elip-d} Ana no leyó el libro en la biblioteca, \textit{y}
          María tampoco.
    \end{xlist}
\end{exe}

Siguiendo a \textcite[278]{ditullio2014}, diferenciamos las instancias de
coordinación con elipsis, como las de \exref{ex:12-elip}, de otros casos en los
que no es necesario plantear una elisión, sino que se puede establecer una
coordinación a nivel sintagmático, como en los ejemplos de
\exref{ex:12-sinelip}:

\begin{exe}
  \ex\label{ex:12-sinelip}
    \begin{xlist}
      \ex\label{ex:12-sinelip-a} Pedro estudia Letras \textit{y} Filosofía.
      \ex\label{ex:12-sinelip-b} Ana leyó el libro en su casa \textit{y} en la
          plaza.
      \ex\label{ex:12-sinelip-c} Ana \textit{y} María son compañeras de trabajo.
    \end{xlist}
\end{exe}

En \exref{ex:12-sinelip-a}, se coordinan dos sintagmas nominales dentro del COD.
En \exref{ex:12-sinelip-b}, los que se coordinan son dos sintagmas
preposicionales, que funcionan como adjuntos dentro del predicado. Por último, en
\exref{ex:12-sinelip-c} se coordinan dos sintagmas nominales, pero dentro del
sujeto y, en este caso en particular, no sería admisible pensar una estructura
alternativa con elipsis, dado que \forma{ser compañeras}, al ser un predicado
simétrico, requiere un sujeto plural. En este sentido, si bien existiría la
posibilidad de plantear versiones completamente redundantes para los casos de
\exref{ex:12-sinelip-a} y \exref{ex:12-sinelip-b}, no existe tal opción para
\exref{ex:12-sinelip-c},\footnote{Si bien es cierto que
\exref{ex:12-sinelip-c} es aceptable si se interpreta que Ana y María son
compañeras de trabajo de una tercera persona, esa lectura no sería simétrica.}
tal como se puede apreciar en los siguientes ejemplos:

\begin{exe}
  \ex
    \begin{xlist}
      \ex Pedro estudia Letras \textit{y (también)} estudia Filosofía.
      \ex Ana leyó el libro en su casa \textit{y (también)} lo leyó en la plaza.
      \ex[*]{Ana es compañera de trabajo \textit{y} María es compañera de
          trabajo.}
    \end{xlist}
\end{exe}

\begin{ejercicio}\label{ej:12-elipsis}
Identifique los elementos que se eliden en el segundo coordinando de las
siguientes oraciones, según se muestra en el ejemplo resuelto~1.

\begin{exe}
  \ex Frodo quiere destruir el anillo y Sam también.
  \ex Aragorn está enamorado de Arwen y Sam, de Rosita.
  \ex Frodo trata bien a Gollum, pero Sam no.
  \ex Boromir no puede ser rey de Gondor, y Faramir tampoco.
  \ex Gandalf no es un mago ambicioso, pero Saruman sí.
  \ex Los elfos desconfían de los hombres y de los enanos.
\end{exe}

\begin{resuelto}
1. Se elide \forma{quiere destruir el anillo}.
\end{resuelto}
\end{ejercicio}

\begin{ejercicio}[Integración]\label{ej:12-meta}
Lea el siguiente fragmento de la nota periodística «¿Puede Meta usar tus fotos
para entrenar su IA?» y resuelva lo pedido respecto de las coordinaciones
señaladas entre corchetes, según se muestra en el ejemplo resuelto~1.

\begin{consignas}
  \item Identifique el tipo de conjunción coordinante.
  \item Analice la estructura del SConj.
  \item Determine si se trata de una coordinación homocategorial o
        heterocategorial.
  \item Considere si hay valores asociados a la conjunción que inciden en la
        interpretación de la coordinación.
\end{consignas}

\begin{textolectura}
(1) La última operación de Meta ha suscitado cierta preocupación entre los
usuarios de [Facebook e Instagram]. La empresa ha anunciado recientemente que
utilizará los datos que estas personas han colgado en sus plataformas para
entrenar su nueva inteligencia artificial (IA), a no ser que expliciten lo
contrario. (2) Un movimiento que ha hecho emerger todo tipo de preguntas sobre
[la protección de datos o su impacto social]. (3) A su vez, ha abierto una
oportunidad para reflexionar sobre la falta de consciencia que tienen los
usuarios [sobre el uso de sus datos en general, pero también sobre las maniobras
normalizadas y ocultas del poder]. [\ldots] (4) (A) [Desde Meta se ha afirmado
que buscan cumplir con (B) [las normativas de privacidad y protección de datos],
pero la desconfianza persiste]. (5) (A) [A su juicio, el problema no gira tanto
en torno a si de aquí en adelante la foto de tu hijo pequeño aparecerá en un
futuro en anuncios de pañales, sino que la cuestión se sitúa en torno a (B) [si
basta con que Meta dé la oportunidad a los usuarios de oponerse a que sus datos
sean utilizados con ese fin o es necesario un consentimiento específico]].
\end{textolectura}
\pendiente{Falta la referencia de la nota (medio, fecha, URL)}

\begin{resuelto}
1. a) Conjunción coordinante copulativa. b) [\etq{SConj/Comp} [\etq{SN/Esp}
Facebook] \textbf{e}\etq{Conj/n} [\etq{SN/Comp} Instagram]]. c) Coordinación
homocategorial entre dos SSNN. d) No hay valores asociados a la conjunción, más
allá del aditivo.
\end{resuelto}
\end{ejercicio}

\begin{recuadro}[Reflexión final]
Como hemos visto en este capítulo, la coordinación se basa en un mecanismo
sumamente productivo, que permite conformar una unidad sintagmática de distintos
grados de complejidad, a partir de la incorporación de una conjunción
coordinante, que se desempeña como núcleo funcional del nuevo sintagma. Mientras
que en la coordinación de sintagmas léxicos (SV, SN, SA, SAdv, SP) el fenómeno se
inscribe dentro de la sintaxis de la oración simple, la coordinación de oraciones
nos lleva a la esfera de las relaciones interoracionales, campo que comparte con
la subordinación.
\end{recuadro}

% --- Secciones finales -------------------------------------------------
\lecturasrecomendadas

\solucionario

\begin{solucion}{\ref{ej:12-predicaciones}}
2. Predicación primaria: \forma{Los transeúntes miraban la escena}. Predicación
secundaria: \forma{Los transeúntes estaban asombrados}. Yuxtaposición:
\forma{Los transeúntes miraban la escena; estaban asombrados}. Coordinación:
\forma{Los transeúntes miraban la escena y estaban asombrados}. Subordinación:
\forma{Los transeúntes, que miraban la escena, estaban asombrados}. 3.
Predicación primaria: \forma{María colgó el cuadro}. Predicación secundaria:
\forma{El cuadro quedó bien centrado}. Yuxtaposición: \forma{María colgó el
cuadro; el cuadro quedó bien centrado}. Coordinación: \forma{María colgó el
cuadro y el cuadro quedó bien centrado}. Subordinación: \forma{María colgó el
cuadro, que quedó bien centrado}.
\end{solucion}

\begin{solucion}{\ref{ej:12-desantis}}
2. Hicimos una colecta para pagarle el legado, sin saber si comprábamos
[\etq{SConj/COD} [\etq{SN/Esp} un tesoro] o\etq{Conj/n} [\etq{SN/Comp} un
montón de diarios viejos]]. 3. [\etq{SConj/OC} [\etq{SF/Esp} Habíamos subido el
baúl por la escalera casi a los saltos], pero\etq{Conj/n} [\etq{SF/Comp} la
visión de aquellos papeles nos paralizó]]. 4. [\etq{SConj/OC} [\etq{SF/Esp}
Balestri sigue siendo un desconocido], y\etq{Conj/n} [\etq{SF/Comp} de su torre
no poseemos más que la sombra de una sombra]].
\end{solucion}

\begin{solucion}{\ref{ej:12-enriquez}}
[\etq{Oración 1} Tenían una madre que rara vez era [\etq{SConj} [\etq{V/Esp}
vista] o\etq{Conj/n} [\etq{V/Comp} convocada]]; la madre Hécate].
[\etq{SConj/Oración 2} [\etq{SF/Esp} En el Enjambre se creía que era una
Leyenda, una mentira], pero\etq{Conj/n} [\etq{SF/Comp} las Luminosas se reían
de tanta ignorancia]]. [\etq{SConj/OS} Y\etq{Conj/n} [\etq{SF/Comp} tenían
hermanas mayores: [\etq{SConj/ApE} [\etq{SN/Esp} Perséfone bajo la tierra, con
flores en el pelo]; y\etq{Conj/n} [\etq{SN/Comp} las Imago, las que se escondían
en lugares [\etq{SConj/M} [\etq{SA/Esp} secos] y\etq{Conj/n} [\etq{SA/Comp}
oscuros]]]]]].\pendiente{La oración que empieza con \textit{Y} se etiqueta como
SConj sin especificador expreso: conviene explicitar el análisis}
\end{solucion}

\begin{solucion}{\ref{ej:12-lecturas}}
2. a) Lectura distributiva: \juicio{*}{\forma{Frodo se lleva bien desde la
infancia y Sam también se lleva bien desde la infancia}}. b) Lectura colectiva:
\forma{Frodo y Sam se llevan bien entre sí desde la infancia}. c) Lectura
ordenada: \juicio{*}{\forma{Frodo y Sam se llevan bien desde la infancia,
respectivamente}}.

3. a) Lectura distributiva: \forma{Frodo viaja a Mordor y Sam también viaja a
Mordor} / \forma{Tanto Frodo como Sam viajan a Mordor (separados)}. b) Lectura
colectiva: \forma{Frodo y Sam viajan a Mordor juntos}. c) Lectura ordenada:
\juicio{*}{\forma{Frodo y Sam viajan a Mordor, respectivamente}}. No es posible
la lectura porque no hay un segundo elemento en el predicado que permita
realizar el ordenamiento.

4. a) Lectura distributiva: se puede armar la secuencia por separado, pero no
mantiene el mismo sentido porque las predicaciones se vuelven reflexivas:
\forma{Sam se odia (a sí mismo) y Gollum también se odia (a sí mismo)} /
\forma{Tanto Sam como Gollum se odian (a sí mismos)}. Al tratarse de un verbo
recíproco, se podría reponer la lectura distributiva si se eliminara el pronombre
\forma{se} y se asignaran las funciones de sujeto y objeto directo
respectivamente a Sam y a Gollum: \forma{Sam odia a Gollum y Gollum odia a Sam}.
b) Lectura colectiva: \forma{Sam y Gollum se odian entre sí / mutuamente}. Es la
única posible, por tratarse de un predicado simétrico, de naturaleza colectiva,
que se predica sobre un grupo. c) Lectura ordenada: no es posible, porque en el
predicado no hay elementos que puedan ordenarse.

5. a) Lectura distributiva: es posible estructuralmente: \forma{Frodo viaja hacia
Mordor y hacia el Bosque de Fangorn, y Aragorn también} / \forma{Tanto Frodo como
Aragorn viajan hacia Mordor y hacia el Bosque de Fangorn}. Podría admitirse tal
lectura si consideramos que primero van a un lugar y después al otro, cada uno
por su cuenta (aunque tal sentido no se corresponde con lo que ocurre en el
relato). b) Lectura colectiva: \forma{Frodo y Aragorn viajan juntos hacia Mordor
y el Bosque de Fangorn}. Es posible si consideramos que van juntos primero a un
lugar y después al otro (aunque tal sentido no se corresponde con lo que ocurre
en el relato). c) Lectura ordenada: \forma{Frodo y Aragorn viajan hacia Mordor y
hacia el Bosque de Fangorn, respectivamente} / \forma{Frodo viaja hacia Mordor y
Aragorn viaja hacia el Bosque de Fangorn}. Este es el sentido compatible con lo
que ocurre en el relato, donde Frodo viaja con Sam hacia el este, mientras que
Aragorn, con otra compañía, viaja hacia el oeste.
\end{solucion}

\begin{solucion}{\ref{ej:12-copulativas}}
2. a) [\etq{SConj/Suj} [\etq{SN/Esp} Meteoritos] y\etq{Conj/n} [\etq{SN/Comp}
volcanes]], en el pasado, provocaron las extinciones masivas documentadas a lo
largo de la historia. b) Coordinación simétrica: es posible cambiar el orden de
los coordinandos. 3. a) El dodo [\etq{SConj/Pred} [\etq{SV/Esp} vivió en la Isla
Mauricio] y\etq{Conj/n} [\etq{SV/Comp} se extinguió hacia el siglo XVII]]. b)
Coordinación asimétrica: se reconoce un valor temporal de sucesión asociado a la
conjunción.\pendiente{El original etiqueta los coordinandos como SF} 4. a) La
compañía [\etq{SConj1/Pred} [\etq{SV/Esp} aislará] y\etq{Conj/n} [\etq{SV/Comp}
cultivará [\etq{SConj2/COD} [\etq{SN/Esp} óvulos] y\etq{Conj/n} [\etq{SN/Comp}
células productoras de esperma]] para luego editar las secuencias de ADN que
coincidan con las del dodo]]. b) SConj1: coordinación asimétrica: se reconoce un
valor temporal de sucesión asociado a la conjunción. SConj2: coordinación
simétrica: es posible cambiar el orden de los coordinandos. 5. a) Estas células
darían lugar a un animal mezcla de [\etq{SConj/Comp} [\etq{SN/Esp} dodo]
y\etq{Conj/n} [\etq{SN/Comp} paloma]]. b) Coordinación simétrica: es posible
cambiar el orden de los coordinandos. 6. a) El mamut lanudo [\etq{SConj1/Pred}
[\etq{SV/Esp} habitó [\etq{SConj2/COD} [\etq{SN/Esp} Asia] y\etq{Conj/n}
[\etq{SN/Comp} Europa]] hace 6 mil años], y\etq{Conj/n} [\etq{SV/Comp} podría
volver a existir gracias a la ciencia]]. b) SConj1: coordinación asimétrica: se
reconoce un valor temporal de sucesión asociado a la conjunción. SConj2:
coordinación simétrica: es posible cambiar el orden de los
coordinandos.\pendiente{El original etiqueta SConj2 como CCL, \textit{Asia} como
SV/Esp y lo califica de «asimétrica: es posible cambiar el orden»}
\end{solucion}

\begin{solucion}{\ref{ej:12-disyuntivas}}
2. a) Los capibaras son dóciles y usualmente los humanos [\etq{SConj/Pred}
[\etq{SV/Esp} los acarician] o\etq{Conj/n} [\etq{SV/Comp} les dan de comer]].
b) Disyunción inclusiva: ambas opciones son posibles. 3. a) Dependiendo de la
subespecie, el periodo de gestación dura [\etq{SConj/CCT} [\etq{SN/Esp} unos 110
días] o\etq{Conj/n} [\etq{SN/Comp} unos 150 días]]. b) Disyunción exclusiva: la
duración del periodo de gestación es distinta para cada subespecie. 4. a) Sus
depredadores naturales más importantes son félidos, como
[\etq{SConj/Comp}\footnote{En (4), consideramos que el núcleo de \forma{como
jaguares o pumas} es \forma{como}, dado que introduce la secuencia que sirve de
ejemplo.} [\etq{SN/Esp} jaguares] o\etq{Conj/n} [\etq{SN/Comp} pumas]]. b)
Disyunción inclusiva: ambas opciones son posibles. 5. a) En la lengua tupí, el
término es \textit{ka'apiûara}, que significa [\etq{SConj/PSO} [\etq{SF/Esp}
\sig{el que come hojas delgadas}] o\etq{Conj/n} [\etq{SN/Comp} \sig{comedor de
hierba}]]. b) Disyunción metalingüística denominativa: introduce una segunda
denominación posible para \textit{ka'apiûara}. 6. a) En 2020, los carpinchos
[\etq{SConj/Pred} [\etq{SV/Esp} invadieron Nordelta] o\etq{Conj/n}, mejor dicho,
[\etq{SV/Comp} volvieron a su territorio autóctono en los humedales del río
Paraná]]. b) Disyunción con valor correctivo: introduce una reformulación que, a
su vez, incorpora la opinión del emisor, que no considera que los carpinchos
hayan invadido la zona, sino que volvieron porque era su territorio autóctono.
\end{solucion}

\begin{solucion}{\ref{ej:12-adversativas}}
2. a) Estos números representan una pérdida ya consolidada del 5 por ciento de
los linajes, es decir, de especies [\etq{SConj/M} [\etq{SA/Esp} distintas],
pero\etq{Conj/n} [\etq{SP/Comp} con un ancestro en común]]. b) La coordinación
adversativa con \forma{pero} establece una contraposición entre los coordinandos,
y el segundo explicita una contraexpectativa: aunque la afirmación del primer
coordinando de que las especies son distintas lleva a inferir que no tienen
ancestros en común, estas especies sí lo tienen (y por eso son del mismo linaje).
3. a) [\etq{SConj/OC} [\etq{SF/Esp} No solo se achicará la biodiversidad de la
fauna], sino que\etq{Conj/n} [\etq{SF/Comp} también se perderán funciones
ecosistémicas cruciales, como la polinización y el control de plagas]]. b) La
coordinación adversativa con \forma{sino} no se interpreta como una oposición
exclusiva, sino como una adición contrastiva, por la presencia del adverbio
\forma{solo} en el primer coordinando y de \forma{también} en el segundo. 4. a)
La paloma migratoria [\etq{SConj/Pred} [\etq{SV/Esp} era la especie de ave más
abundante en Norteamérica], pero\etq{Conj/n} [\etq{SV/Comp} ahora está
completamente extinta]]. b) La coordinación adversativa con \forma{pero} se
interpreta como valorizadora del segundo coordinando. Los coordinandos se
relacionan en una jerarquía de valores (antes era abundante, ahora está extinta),
y se prioriza el segundo. El orden de los coordinandos podría cambiarse
(\forma{La paloma migratoria ahora está completamente extinta, pero era la
especie más abundante en Norteamérica}); en ese caso, cambiaría también la
lectura valorizadora, dado que siempre se le da prominencia al segundo
elemento.\pendiente{En el original, las respuestas 3 y 4 están intercambiadas} 5.
a) [\etq{SConj/OC} [\etq{SF/Esp} El ser humano no se siente parte del
ecosistema], sino que\etq{Conj/n} [\etq{SF/Comp} está totalmente desconectado de
la naturaleza de la que forma parte]]. b) La coordinación adversativa con
\forma{sino} se interpreta como oposición exclusiva, en la que el segundo
coordinando reemplaza al elemento negado en primer lugar. 6. a) [\etq{SConj/OC}
[\etq{SF/Esp} Hay medidas para frenar este conflicto], pero\etq{Conj/n}
[\etq{SF/Comp} se tienen que tomar muy rápido]]. b) La coordinación adversativa
con \forma{pero} establece una contraposición entre los coordinandos, aunque, en
realidad, la contraposición se da entre las inferencias que se realizan sobre
cada coordinando: hay medidas para frenar el conflicto, de lo cual se infiere
cierto optimismo, pero las medidas se deben tomar muy rápido, de lo cual se
infiere que el panorama optimista requiere algo más que posibles medidas. Esta
interpretación se valida en una afirmación que hace el autor unas líneas más
adelante, según la cual el optimismo debe alimentarse de opciones concretas.
\end{solucion}

\begin{solucion}{\ref{ej:12-elipsis}}
2. Se elide \forma{está enamorado}. 3. Se elide \forma{trata bien a Gollum}. 4.
Se elide \forma{puede ser rey de Gondor}. 5. Se elide \forma{es un mago
ambicioso}. 6. No hay elipsis, sino una coordinación copulativa de dos SSPP. El
SConj funciona como el complemento preposicional de régimen del verbo
\forma{desconfiar}.
\end{solucion}

\begin{solucion}{\ref{ej:12-meta}}
2. a) Conjunción coordinante disyuntiva \forma{o}. b) [\etq{SConj/Comp}
[\etq{SN/Esp} la protección de datos] \textbf{o}\etq{Conj/n} [\etq{SN/Comp} su
impacto social]]. c) Coordinación homocategorial entre dos SSNN. d) La
disyunción se interpreta como inclusiva: no se impone la elección de una u otra
opción.

3. a) Conjunción coordinante adversativa \forma{pero}. b) [\etq{SConj/CPR}
[\etq{SP/Esp} sobre el uso de sus datos en general], \textbf{pero}\etq{Conj/n}
también [\etq{SP/Comp} sobre las maniobras normalizadas y ocultas del poder]].
c) Coordinación homocategorial entre dos SSPP. d) La conjunción está acompañada
por el adverbio \forma{también}, que funciona como marcador discursivo. La
coordinación se interpreta con una lectura valorizadora, en la que se da
prominencia al segundo constituyente.\pendiente{El original incluye en el primer
coordinando \textit{sobre la falta de consciencia que tienen los usuarios}}

4. (A) a) Conjunción coordinante adversativa \forma{pero}. b) [\etq{SConj/OC}
[\etq{SF/Esp} Desde Meta se ha afirmado que buscan cumplir con las normativas de
privacidad y protección de datos], \textbf{pero}\etq{Conj/n} [\etq{SF/Comp} la
desconfianza persiste]]. c) Coordinación homocategorial entre dos oraciones. d)
Los dos coordinandos son sintagmas de flexión, por lo que el SConj se constituye
como una oración compuesta. Se interpreta un valor de contraexpectativa porque
Meta afirma que cumple con la normativa pertinente y se infiere que eso debería
generar confianza en los usuarios, pero eso no sucede. (B) a) Conjunción
coordinante copulativa \forma{y}. b) [\etq{SConj/Comp} [\etq{SN/Esp} privacidad]
\textbf{y}\etq{Conj/n} [\etq{SN/Comp} protección de datos]]. c) Coordinación
homocategorial entre dos SSNN. d) No hay valores asociados a la conjunción.

5. (A) a) Conjunción coordinante adversativa \forma{sino}. b) [\etq{SConj/OC}
[\etq{SF/Esp} A su juicio, el problema no gira tanto en torno a si de aquí en
adelante la foto de tu hijo pequeño aparecerá en un futuro en anuncios de
pañales], \textbf{sino que}\etq{Conj/n} [\etq{SF/Comp} la cuestión se sitúa en
torno a si basta con que Meta dé la oportunidad a los usuarios de oponerse a que
sus datos sean utilizados con ese fin o es necesario un consentimiento
específico]]. c) Coordinación homocategorial entre dos oraciones. d) La
coordinación se interpreta con valor de oposición exclusiva, en la que el
segundo elemento reemplaza al primero. (B) a) Conjunción coordinante disyuntiva
\forma{o}. b) [\etq{SConj} [\etq{SF/Esp} si basta con que Meta dé la
oportunidad a los usuarios de oponerse a que sus datos sean utilizados con ese
fin] \textbf{o}\etq{Conj/n} [\etq{SF/Comp} (si) es necesario un consentimiento
específico]]. c) Coordinación homocategorial entre dos oraciones. d) La
disyunción se interpreta como exclusiva: se impone la elección de una u otra
opción. En el segundo coordinando hay una elipsis del subordinante \forma{si},
que ya está mencionado en el primer constituyente.\pendiente{El original
etiqueta este SConj como OComp, pero está incluido como término de \textit{en
torno a}}
\end{solucion}
